\documentclass[14pt,a4paper]{extarticle}

\usepackage{cmap}
\usepackage[T2A]{fontenc}
\usepackage[utf8]{inputenc}
\usepackage[russian]{babel}
\usepackage{newtxtext,newtxmath}

\usepackage{amsmath}
\usepackage{geometry}
\usepackage{setspace}
\usepackage{indentfirst}
\usepackage{graphicx}
\usepackage{array}
\usepackage{caption}
\usepackage{booktabs}
\usepackage{siunitx}
\sisetup{output-decimal-marker={,}, separate-uncertainty, per-mode=symbol}
\DeclareSIUnit\cm{\text{см}}
\captionsetup[table]{labelsep=endash, justification=raggedright, singlelinecheck=false}
\captionsetup[figure]{labelsep=endash}
\usepackage{pgfplots}
\pgfplotsset{compat=1.18}

\geometry{left=30mm, right=10mm, top=20mm, bottom=20mm}
\onehalfspacing
\setlength{\parindent}{1.25cm}
\pagestyle{plain}

\begin{document}

\begin{titlepage}
\begin{center}
\onehalfspacing
Белорусский государственный университет\\
Факультет радиофизики и компьютерных технологий

\vspace{4cm}

\textbf{ОТЧЁТ}\\
по лабораторной работе №1\\[0.3cm]
\textbf{«Изучение линз и оптических систем»}

\vspace{4cm}

\begin{flushright}
\begin{tabular}{@{}p{7.5cm}@{}}
\raggedright
Выполнил:\\
студент 2 курса 8 группы\\
Нацевский М. А.\\
Преподаватель:\\
Трухан С. В.\\
\end{tabular}
\end{flushright}

\vfill
Минск, 2026
\end{center}
\end{titlepage}
\setcounter{page}{2}

\section*{Цель работы}

Изучить методы расчёта центрированных оптических систем. Измерить фокусные расстояния тонкой положительной и тонкой отрицательной линз.

\section*{Выполнение работы}

\textbf{Линза №1, способ 1.} На оптической скамье получали чёткое изображение освещённой шкалы на экране и измеряли расстояния от линзы до предмета $a_1$ и до изображения $a_2$. Результаты приведены в табл.~\ref{tab:lens1}, где $L = a_1 + a_2$ --- расстояние от предмета до экрана.

\begin{table}[htbp]
\caption{Результаты измерений для линзы №1 (способ 1)}
\label{tab:lens1}
\centering
\small
\setlength{\tabcolsep}{4pt}
\begin{tabular}{
S[table-format=2.0]
S[table-format=2.2]
S[table-format=2.2]
S[table-format=3.0]
S[table-format=2.3]
S[table-format=2.3]
S[table-format=2.3]
S[table-format=2.2]
}
\toprule
{№} & {$a_1$, см} & {$a_2$, см} & {$L$, см} & {$\dfrac{1}{a_1}$, $10^{-3}$ см$^{-1}$} & {$\dfrac{1}{a_2}$, $10^{-3}$ см$^{-1}$} & {$\dfrac{1}{a_1}+\dfrac{1}{a_2}$, $10^{-3}$ см$^{-1}$} & {$f_i$, см} \\
\midrule
1 & 43 & 60 & 103 & 23,256 & 16,667 & 39,922 & 25,05 \\
2 & 60,1 & 42,9 & 103 & 16,639 & 23,310 & 39,949 & 25,03 \\
3 & 40,1 & 65,9 & 106 & 24,938 & 15,175 & 40,112 & 24,93 \\
4 & 65 & 41 & 106 & 15,385 & 24,390 & 39,775 & 25,14 \\
5 & 39 & 70 & 109 & 25,641 & 14,286 & 39,927 & 25,05 \\
6 & 70 & 39 & 109 & 14,286 & 25,641 & 39,927 & 25,05 \\
7 & 70,55 & 41,45 & 112 & 14,174 & 24,125 & 38,300 & 26,11 \\
8 & 37,5 & 74,5 & 112 & 26,667 & 13,423 & 40,089 & 24,94 \\
9 & 78 & 37 & 115 & 12,821 & 27,027 & 39,848 & 25,10 \\
10 & 36,5 & 78,5 & 115 & 27,397 & 12,739 & 40,136 & 24,92 \\
\bottomrule
\end{tabular}
\end{table}

\textbf{Линза №1, способ 2.} При неизменном расстоянии между предметом и экраном $L = \SI{103}{\cm}$ найдены два положения линзы, дающие чёткое изображение: \SI{60}{\cm} и \SI{44}{\cm}. Расстояние между ними $l = \SI{16}{\cm}$.

\textbf{Линза №1, способ 3.} С помощью зрительной трубы, установленной на бесконечность, найдено положение линзы, при котором предмет находится в её фокальной плоскости. Расстояние от предмета до линзы составило $f \approx \SI{25}{\cm}$.

\textbf{Линза №2, способ 1.} В качестве вспомогательной использовалась линза №1. Изображение, создаваемое ею, находилось в точке $S_1 = \SI{105}{\cm}$. При отрицательной линзе, установленной в точке $x_p = \SI{95}{\cm}$, изображение сместилось в точку $S_2 = \SI{115}{\cm}$.

\textbf{Линза №2, способ 2.} Отрицательная линза перемещалась до появления чёткого изображения в зрительной трубе, установленной на бесконечность. Это произошло при $x_p = \SI{85}{\cm}$ и $S_1 = \SI{105}{\cm}$.

Измерения для толстой линзы №3 не проводились, поэтому положения её фокальных и главных плоскостей в отчёте не определяются.

\section*{Обработка экспериментальных данных}

\textbf{Линза №1, способ 1.} Формула тонкой линзы для действительного предмета имеет вид
\[
\frac{1}{a_2} + \frac{1}{a_1} = \frac{1}{f},
\qquad\text{или}\qquad
y = \beta_0 + \beta_1 x,
\]
где $y = 1/a_2$, $x = 1/a_1$, $\beta_0 = 1/f$, $\beta_1 = -1$.

Значение в измерении №7 заметно отличается от остальных. Проверка по критерию Граббса для величины $1/a_1 + 1/a_2$ даёт $G = 2{,}78$, что больше критического значения $G_{\text{кр}} = 2{,}29$ при $n = 10$ и уровне значимости $0{,}05$. Измерение №7 признано промахом и исключено из расчёта.

По оставшимся $N = 9$ измерениям при $\beta_1 = -1$:
\[
\beta_0 = \frac{1}{N}\sum_{n=1}^{N}\left(y_n + x_n\right) = \SI{39,965e-3}{\cm^{-1}},
\qquad
s_{\beta_0} = \SI{0,041e-3}{\cm^{-1}}.
\]
Доверительный интервал при вероятности $0{,}95$ (коэффициент Стьюдента $t = 2{,}31$) равен $\Delta\beta_0 = t\,s_{\beta_0} = \SI{0,095e-3}{\cm^{-1}}$. Отсюда
\[
f = \frac{1}{\beta_0} = \SI{25,02}{\cm},
\qquad
\Delta f = f^2 \Delta\beta_0 = \SI{0,06}{\cm}.
\]

Зависимость $1/a_2 = F(1/a_1)$ приведена на рис.~\ref{fig:lens1}. Для проверки выполнена аппроксимация методом наименьших квадратов со свободным угловым коэффициентом: $\beta_1 = -0{,}984 \pm 0{,}012$, $f = (25{,}23 \pm 0{,}16)$~см. Угловой коэффициент близок к теоретическому значению $-1$.

\begin{figure}[htbp]
\centering
\begin{tikzpicture}
\begin{axis}[
  width=0.9\textwidth, height=8cm,
  xlabel={$1/a_1$, $10^{-3}$ см$^{-1}$}, ylabel={$1/a_2$, $10^{-3}$ см$^{-1}$},
  grid=major, xmin=11, xmax=29, ymin=10, ymax=30,
  legend pos=north east, legend cell align=left,
  /pgf/number format/use comma
]
\addplot[only marks, mark=*, blue] coordinates {
(23.256,16.667) (16.639,23.310) (24.938,15.175) (15.385,24.390) (25.641,14.286)
(14.286,25.641) (26.667,13.423) (12.821,27.027) (27.397,12.739)
};
\addplot[only marks, mark=o, mark size=3pt, thick, black] coordinates {(14.174,24.125)};
\addplot[red, thick, domain=11:29] {39.965 - x};
\legend{Эксперимент, Промах (№7), {$y = \beta_0 - x$}}
\end{axis}
\end{tikzpicture}
\caption{Зависимость $1/a_2$ от $1/a_1$ для линзы №1}
\label{fig:lens1}
\end{figure}

\textbf{Линза №1, способ 2.} Фокусное расстояние:
\[
f = \frac{L^2 - l^2}{4L} = \frac{103^2 - 16^2}{4 \cdot 103}~\text{см} = \SI{25,13}{\cm}.
\]
Погрешность отсчёта положения по линейке принята равной $\delta = \SI{0,5}{\cm}$, тогда $\Delta L = \Delta l = \sqrt{2}\,\delta \approx \SI{0,7}{\cm}$ и
\[
\Delta f = \sqrt{\left(\frac{L^2 + l^2}{4L^2}\,\Delta L\right)^2 + \left(\frac{l}{2L}\,\Delta l\right)^2} \approx \SI{0,2}{\cm}.
\]

\textbf{Линза №1, способ 3.} Фокусное расстояние определено однократно, погрешность наводки на резкость принята равной \SI{1}{\cm}: $f = (25 \pm 1)$~см.

\textbf{Линза №2, способ 1.} Мнимый источник и изображение лежат справа от линзы, поэтому
\[
a_1 = S_1 - x_p = \SI{10}{\cm},
\qquad
a_2 = S_2 - x_p = \SI{20}{\cm}.
\]
Из формулы линзы $1/a_2 - 1/a_1 = 1/f$ получаем
\[
f = \frac{a_1 a_2}{a_1 - a_2} = \frac{10 \cdot 20}{10 - 20}~\text{см} = \SI{-20}{\cm}.
\]
При $\delta = \SI{0,5}{\cm}$ для каждого из отсчётов $S_1$, $S_2$, $x_p$ погрешность составляет
\[
\Delta f = \frac{\delta\sqrt{a_2^4 + a_1^4 + \left(a_2^2 - a_1^2\right)^2}}{\left(a_1 - a_2\right)^2} \approx \SI{3}{\cm}.
\]

\textbf{Линза №2, способ 2.} Мнимый источник совпадает с фокусом линзы, а фокусное расстояние отрицательно:
\[
|f| = S_1 - x_p = 105 - 85 = \SI{20}{\cm},
\qquad
\Delta f = \sqrt{2}\,\delta \approx \SI{1}{\cm}.
\]

Результаты всех измерений сведены в табл.~\ref{tab:results}.

\begin{table}[htbp]
\caption{Фокусные расстояния линз}
\label{tab:results}
\centering
\begin{tabular}{lll}
\toprule
Линза & Способ & $f$, см \\
\midrule
№1 (положительная) & 1: формула линзы, МНК & $25{,}02 \pm 0{,}06$ \\
№1 (положительная) & 2: метод смещения линзы & $25{,}1 \pm 0{,}2$ \\
№1 (положительная) & 3: зрительная труба & $25 \pm 1$ \\
№2 (отрицательная) & 1: формула линзы & $-20 \pm 3$ \\
№2 (отрицательная) & 2: зрительная труба & $-20 \pm 1$ \\
\bottomrule
\end{tabular}
\end{table}

\section*{Вывод}

Изучены методы расчёта центрированных оптических систем и измерены фокусные расстояния тонких линз. Для положительной линзы №1 три независимых способа дали согласующиеся в пределах погрешности значения; наиболее точный результат получен по формуле линзы: $f = (25{,}02 \pm 0{,}06)$~см. Для отрицательной линзы №2 оба способа дали одинаковое значение, наиболее точный результат получен с помощью зрительной трубы: $f = (-20 \pm 1)$~см.

\end{document}
